% Generated by Listas/make_statements.py from pratica-lecture-05.md, item 3. Do not edit.
\textbf{Source.} Muñoz-Garcia \& Toro-Gonzalez (2019), \emph{Strategy and Game Theory: Practice Exercises with Answers}, Ch. 4, Exercise 18 ``Backward Induction in Wage Negotiations'', pp.~200--201 {[}PDF 212--213{]}. Answer in the source: yes, a worked answer follows the exercise. Reproduced verbatim (statement only; the book's worked answer follows it).



\begin{quote}

Consider the following bargaining game between a labor union and a firm. In the first stage, the labor union chooses the wage level, \(w\), that all workers will receive. In the second stage, the firm responds to the wage \(w\) by choosing the number of workers it hires, \(h \in [0, 1]\). The labor union's payoff function is \(u_L(w, h) = w \cdot h\), thus being increasing in wages and in the number of workers hired. The firm's profit is



\[\pi(w, h) = \left(h - \frac{h^2}{2}\right) - wh\]



Intuitively, revenue is increasing in the number of workers hired, \(h\), but at a decreasing rate (i.e., \(h - \frac{h^2}{2}\) is concave in \(h\)), reaching a maximum when the firm hires all workers, i.e., \(h = 1\) where revenue becomes 1/2.



\textbf{Part (a)} Applying backward induction, analyze the optimal strategy of the last mover (firm). Find the firm's best response function \(h(w)\).



\textbf{Part (b)} Let us now move on to the first mover in the game (labor union). Anticipating the best response function of the firm that you found in part (a), \(h(w)\), determine the labor union's optimal wage, \(w^*\).

\end{quote}

